% !TeX root = ../Book1.tex
\begin{tikzpicture}[scale=0.9]
	\begin{scope}
		\coordinate (p1) at (0,0);
		\coordinate (p2) at (3,0);
		\coordinate (p3) at (2.5,2.1);
		\coordinate (p4) at (-0.3,1.7);
		\coordinate (z) at (1.14043,0.95796);
		\draw[gray!55] (p1)--(p2)--(p3)--(p4)--cycle;
		\draw[very thick,blue!70!black] (p1)--(p3);
		\draw[very thick,orange!85!black] (p2)--(p4);
		\fill[blue!70!black] (p1) circle(2pt) node[below left] {\(p_1\)};
		\fill[blue!70!black] (p3) circle(2pt) node[above right] {\(p_3\)};
		\fill[orange!85!black] (p2) circle(2pt) node[below right] {\(p_2\)};
		\fill[orange!85!black] (p4) circle(2pt) node[above left] {\(p_4\)};
		\fill (z) circle(2pt) node[above left] {\(z\)};
		\node at (1.3,-0.7) {两条线段相交};
	\end{scope}
	\begin{scope}[xshift=6cm]
		\coordinate (q1) at (0,0);
		\coordinate (q2) at (3,0);
		\coordinate (q3) at (0.8,2.1);
		\coordinate (q4) at (1.1,0.75);
		\filldraw[draw=blue!70!black,fill=blue!8,thick] (q1)--(q2)--(q3)--cycle;
		\fill[blue!70!black] (q1) circle(2pt) node[below left] {\(p_1\)};
		\fill[blue!70!black] (q2) circle(2pt) node[below right] {\(p_2\)};
		\fill[blue!70!black] (q3) circle(2pt) node[above] {\(p_3\)};
		\fill[orange!85!black] (q4) circle(2.5pt) node[right] {\(p_4=z\)};
		\node at (1.3,-0.7) {一点位于三角形内};
	\end{scope}
\end{tikzpicture}
